\begin{thebibliography}{99}
\footnotesize
\bibitem{detrano1989} R.~Detrano \emph{et al.}, ``International application of a new probability algorithm for the diagnosis of coronary artery disease,'' \emph{Am. J. Cardiol.}, vol.~64, no.~5, pp.~304--310, 1989.
\bibitem{arabasadi2017} Z.~Arabasadi \emph{et al.}, ``Computer aided decision making for heart disease detection using hybrid neural network-genetic algorithm,'' \emph{Comput. Methods Programs Biomed.}, vol.~141, pp.~19--26, 2017.
\bibitem{haq2018} A.~U. Haq \emph{et al.}, ``A hybrid intelligent system framework for the prediction of heart disease using machine learning algorithms,'' \emph{Mobile Inf. Syst.}, vol.~2018, pp.~1--21, 2018.
\bibitem{mohan2019} S.~Mohan \emph{et al.}, ``Effective heart disease prediction using hybrid machine learning techniques,'' \emph{IEEE Access}, vol.~7, pp.~81542--81554, 2019.
\bibitem{khan2021} M.~A. Khan, ``An IoT framework for heart disease prediction based on MDCNN classifier,'' \emph{IEEE Access}, vol.~9, pp.~34717--34727, 2021.
\bibitem{wagner2020} P.~Wagner \emph{et al.}, ``PTB-XL, a large publicly available electrocardiography dataset,'' \emph{Sci. Data}, vol.~7, no.~1, p.~154, 2020.
\bibitem{latha2019} C.~B.~C. Latha and S.~C. Jeeva, ``Improving the accuracy of prediction of heart disease risk based on ensemble classification techniques,'' \emph{Inform. Med. Unlocked}, vol.~16, p.~100203, 2019.
\bibitem{gao2021} X.-Y. Gao \emph{et al.}, ``Improving the accuracy for analyzing heart diseases prediction based on the ensemble method,'' \emph{Complexity}, vol.~2021, p.~6663455, 2021.
\bibitem{hannun2019} A.~Y. Hannun \emph{et al.}, ``Cardiologist-level arrhythmia detection and classification in ambulatory electrocardiograms using a deep neural network,'' \emph{Nat. Med.}, vol.~25, no.~1, pp.~65--69, 2019.
\bibitem{christ2018} M.~Christ \emph{et al.}, ``Time series feature extraction on basis of scalable hypothesis tests (tsfresh),'' \emph{Neurocomputing}, vol.~307, pp.~72--77, 2018.
\bibitem{litjens2017} G.~Litjens \emph{et al.}, ``A survey on deep learning in medical image analysis,'' \emph{Med. Image Anal.}, vol.~42, pp.~60--88, 2017.
\bibitem{attia2019} Z.~I. Attia \emph{et al.}, ``Screening for cardiac contractile dysfunction using an artificial intelligence-enabled electrocardiogram,'' \emph{Nat. Med.}, vol.~25, no.~1, pp.~70--74, 2019.
\bibitem{zhang2022} X.~Zhang \emph{et al.}, ``Multi-view learning framework combining ECG and clinical data for cardiac diagnosis,'' \emph{IEEE J. Biomed. Health Inform.}, vol.~26, no.~3, pp.~1012--1020, 2022.
\bibitem{vaswani2017} A.~Vaswani \emph{et al.}, ``Attention is all you need,'' in \emph{Adv. Neural Inf. Process. Syst.}, 2017, pp.~5998--6008.
\bibitem{lundberg2017} S.~M. Lundberg and S.-I. Lee, ``A unified approach to interpreting model predictions,'' in \emph{Adv. Neural Inf. Process. Syst.}, 2017, pp.~4765--4774.
\bibitem{lundberg2018tree} S.~M. Lundberg \emph{et al.}, ``Consistent individualized feature attribution for tree ensembles,'' arXiv:1802.03888, 2018.
\bibitem{chawla2002} N.~V. Chawla \emph{et al.}, ``SMOTE: Synthetic minority over-sampling technique,'' \emph{J. Artif. Intell. Res.}, vol.~16, pp.~321--357, 2002.
\bibitem{alizadehsani2013} R.~Alizadehsani \emph{et al.}, ``A data mining approach for diagnosis of coronary artery disease,'' \emph{Comput. Methods Programs Biomed.}, vol.~111, no.~1, pp.~52--61, 2013.
\bibitem{lyu2020} Q.~Lyu \emph{et al.}, ``Multi-label cardiac arrhythmia classification from electrocardiogram with deep learning,'' in \emph{Proc. ICASSP}, 2020, pp.~1104--1108.
\bibitem{sharma2020} H.~Sharma and N.~Rizvi, ``Prediction of heart disease using machine learning algorithms,'' \emph{Int. J. Inf. Technol.}, vol.~12, no.~2, pp.~355--363, 2020.
\bibitem{pouriyeh2017} S.~Pouriyeh \emph{et al.}, ``A comprehensive investigation and comparison of machine learning techniques in the domain of heart disease classification,'' in \emph{Proc. IEEE ISCC}, 2017, pp.~204--207.
\bibitem{ribeiro2020} A.~H. Ribeiro \emph{et al.}, ``Automatic diagnosis of the 12-lead ECG using a deep neural network,'' \emph{Nat. Commun.}, vol.~11, no.~1, p.~1760, 2020.
\bibitem{lundberg2018sepsis} S.~M. Lundberg \emph{et al.}, ``Explainable machine-learning predictions for the prevention of hypoxaemia during surgery,'' \emph{Nat. Biomed. Eng.}, vol.~2, no.~10, pp.~749--760, 2018.
\bibitem{bi2021} W.~Bi \emph{et al.}, ``Application of support vector machine and unsupervised learning for detecting COVID-19 severity using cardiac MRI-derived radiomics features,'' \emph{Diagnostics}, vol.~11, no.~3, p.~545, 2021.
\bibitem{benjamin2019} E.~J. Benjamin \emph{et al.}, ``Heart disease and stroke statistics -- 2019 update,'' \emph{Circulation}, vol.~139, no.~10, pp.~e56--e528, 2019.
\bibitem{tison2018} G.~H. Tison \emph{et al.}, ``Passive detection of atrial fibrillation using a commercially available smartwatch,'' \emph{JAMA Cardiol.}, vol.~3, no.~5, pp.~409--416, 2018.
\bibitem{rossi2021} R.~A. Rossi \emph{et al.}, ``Precision medicine in patients with heart failure,'' \emph{Heart Fail. Clin.}, vol.~17, no.~2, pp.~199--214, 2021.
\bibitem{chen2016} T.~Chen and C.~Guestrin, ``XGBoost: A scalable tree boosting system,'' in \emph{Proc. ACM SIGKDD}, 2016, pp.~785--794.
\bibitem{breiman2001} L.~Breiman, ``Random forests,'' \emph{Mach. Learn.}, vol.~45, no.~1, pp.~5--32, 2001.
\bibitem{cortes1995} C.~Cortes and V.~Vapnik, ``Support-vector networks,'' \emph{Mach. Learn.}, vol.~20, no.~3, pp.~273--297, 1995.
\bibitem{kingma2015} D.~P. Kingma and J.~Ba, ``Adam: A method for stochastic optimization,'' in \emph{Proc. ICLR}, 2015.
\bibitem{ioffe2015} S.~Ioffe and C.~Szegedy, ``Batch normalization,'' in \emph{Proc. ICML}, 2015, pp.~448--456.
\bibitem{srivastava2014} N.~Srivastava \emph{et al.}, ``Dropout: A simple way to prevent neural networks from overfitting,'' \emph{J. Mach. Learn. Res.}, vol.~15, no.~1, pp.~1929--1958, 2014.
\bibitem{pedregosa2011} F.~Pedregosa \emph{et al.}, ``Scikit-learn: Machine learning in Python,'' \emph{J. Mach. Learn. Res.}, vol.~12, pp.~2825--2830, 2011.
\bibitem{james2021} G.~James \emph{et al.}, \emph{An Introduction to Statistical Learning with Applications in R}, 2nd~ed. New York, NY, USA: Springer, 2021.
\bibitem{who2023} World Health Organization, ``Cardiovascular diseases (CVDs) -- key facts,'' WHO Fact Sheet, 2023. [Online]. Available: \url{https://www.who.int/news-room/fact-sheets/detail/cardiovascular-diseases-(cvds)}
\bibitem{hossain2023} M.~M. Hossain, M.~R. Islam, S.~A. Hossain, and M.~A. Moni, ``Heartbeat classification based on single Lead-II ECG using deep learning,'' \emph{Healthcare}, vol.~11, no.~15, p.~2183, 2023.
\bibitem{janosi1989} A.~Janosi, W.~Steinbrunn, M.~Pfisterer, and R.~Detrano, ``Heart Disease,'' UCI Machine Learning Repository, 1989. doi: 10.24432/C52P4X.
\bibitem{nadeau2003} C.~Nadeau and Y.~Bengio, ``Inference for the generalization error,'' \emph{Mach. Learn.}, vol.~52, no.~3, pp.~239--281, 2003.
\bibitem{bouckaert2004} R.~R. Bouckaert and E.~Frank, ``Evaluating the replicability of significance tests for comparing learning algorithms,'' in \emph{Proc. PAKDD}, 2004, pp.~3--12.
\bibitem{holm1979} S.~Holm, ``A simple sequentially rejective multiple test procedure,'' \emph{Scand. J. Statist.}, vol.~6, no.~2, pp.~65--70, 1979.
\bibitem{ojala2010} M.~Ojala and G.~C. Garriga, ``Permutation tests for studying classifier performance,'' \emph{J. Mach. Learn. Res.}, vol.~11, pp.~1833--1863, 2010.
\bibitem{vancalster2019} B.~Van~Calster \emph{et al.}, ``Calibration: the Achilles heel of predictive analytics,'' \emph{BMC Med.}, vol.~17, p.~230, 2019.
\bibitem{huang2020fusion} S.-C. Huang \emph{et al.}, ``Fusion of medical imaging and electronic health records using deep learning: a systematic review and implementation guidelines,'' \emph{npj Digit. Med.}, vol.~3, p.~136, 2020.
\bibitem{kumar2020problems} I.~E. Kumar \emph{et al.}, ``Problems with Shapley-value-based explanations as feature importance measures,'' in \emph{Proc. ICML}, 2020, pp.~5491--5500.
\bibitem{molnar2022pitfalls} C.~Molnar \emph{et al.}, ``General pitfalls of model-agnostic interpretation methods for machine learning models,'' in \emph{xxAI -- Beyond Explainable AI}, LNCS vol.~13200, Springer, 2022, pp.~39--68.
\bibitem{slack2020fooling} D.~Slack \emph{et al.}, ``Fooling LIME and SHAP: Adversarial attacks on post hoc explanation methods,'' in \emph{Proc. AAAI/ACM Conf. AI, Ethics, Society}, 2020, pp.~180--186.
\bibitem{he2008adasyn} H.~He \emph{et al.}, ``ADASYN: Adaptive synthetic sampling approach for imbalanced learning,'' in \emph{Proc. IEEE IJCNN}, 2008, pp.~1322--1328.
\bibitem{han2005borderline} H.~Han, W.-Y. Wang, and B.-H. Mao, ``Borderline-SMOTE: A new over-sampling method in imbalanced data sets learning,'' in \emph{Proc. ICIC}, 2005, pp.~878--887.
\bibitem{gorishniy2021} Y.~Gorishniy \emph{et al.}, ``Revisiting deep learning models for tabular data,'' in \emph{Adv. Neural Inf. Process. Syst.}, 2021, pp.~18932--18943.
\bibitem{vabalas2019} A.~Vabalas \emph{et al.}, ``Machine learning algorithm validation with a limited sample size,'' \emph{PLoS ONE}, vol.~14, no.~11, p.~e0224365, 2019.
\end{thebibliography}
